\documentclass[11pt,a4paper]{article}

% Compile with: latexmk -pdf gw_speed_monotonicity.tex
% The accompanying .latexmkrc puts all build products in build/.
\usepackage[T1]{fontenc}
\usepackage[utf8]{inputenc}
\usepackage{amsmath,amsthm,mathtools}
\usepackage{newpxtext,newpxmath}
\usepackage{microtype}
\usepackage[a4paper,top=25mm,bottom=26mm,left=29mm,right=29mm,headsep=7mm]{geometry}
\usepackage{xcolor}
\usepackage{titlesec}
\usepackage{fancyhdr}
\usepackage{aliascnt}
\usepackage[unicode,hidelinks]{hyperref}
\usepackage[nameinlink,noabbrev]{cleveref}

\definecolor{ink}{gray}{0.12}
\definecolor{muted}{gray}{0.40}
\color{ink}
\hypersetup{
  pdftitle={A fixed-point correction for speed monotonicity},
  pdfsubject={Galton--Watson trees with at least two offspring},
  pdfauthor={},
  pdfkeywords={Galton--Watson tree, biased random walk, speed monotonicity, recursive distributional equation, fixed-point correction}
}
\titleformat{\section}{\large\bfseries}{\thesection}{0.7em}{}
\titleformat{\subsection}{\normalsize\bfseries}{\thesubsection}{0.6em}{}
\titlespacing*{\section}{0pt}{2.1ex plus .6ex minus .2ex}{1ex}
\titlespacing*{\subsection}{0pt}{1.8ex plus .5ex minus .2ex}{.7ex}
\setlength{\parindent}{1.15em}
\setlength{\parskip}{0.22em}
\linespread{1.035}
\setlength{\emergencystretch}{1.8em}
\setlength{\headheight}{14pt}
\pagestyle{fancy}
\fancyhf{}
\fancyhead[L]{\footnotesize\color{muted}A fixed-point correction for speed monotonicity}
\fancyhead[R]{\footnotesize\color{muted}Research note}
\fancyfoot[C]{\footnotesize\thepage}
\renewcommand{\headrulewidth}{0.25pt}
\renewcommand{\footrulewidth}{0pt}
\fancypagestyle{firstpage}{\fancyhf{}\fancyfoot[C]{\footnotesize\thepage}\renewcommand{\headrulewidth}{0pt}}
\numberwithin{equation}{section}
\allowdisplaybreaks[1]

\newtheoremstyle{notethm}{.9em}{.8em}{\itshape}{} {\bfseries}{.}{.5em}{}
\theoremstyle{notethm}
\newtheorem{theorem}{Theorem}[section]
\newaliascnt{proposition}{theorem}
\newtheorem{proposition}[proposition]{Proposition}
\aliascntresetthe{proposition}
\newaliascnt{lemma}{theorem}
\newtheorem{lemma}[lemma]{Lemma}
\aliascntresetthe{lemma}
\newaliascnt{corollary}{theorem}
\newtheorem{corollary}[corollary]{Corollary}
\aliascntresetthe{corollary}
\theoremstyle{definition}
\newaliascnt{remark}{theorem}
\newtheorem{remark}[remark]{Remark}
\aliascntresetthe{remark}
\renewcommand{\qedsymbol}{\ensuremath{\square}}
\crefname{equation}{equation}{equations}
\crefname{lemma}{Lemma}{Lemmas}
\crefname{proposition}{Proposition}{Propositions}
\crefname{theorem}{Theorem}{Theorems}
\crefname{corollary}{Corollary}{Corollaries}
\crefname{remark}{Remark}{Remarks}
\crefname{section}{Section}{Sections}

\newcommand{\E}{\mathbb E}
\newcommand{\PP}{\mathbb P}
\newcommand{\R}{\mathbb R}
\newcommand{\Law}{\operatorname{Law}}
\newcommand{\Cov}{\operatorname{Cov}}
\newcommand{\id}{\mathrm{id}}
\newcommand{\dd}{\,\mathrm d}
\newcommand{\st}{\leq_{\mathrm{st}}}
\newcommand{\Dplus}{D^{+}}
\newcommand{\cJ}{\mathcal J}
\newcommand{\cL}{\mathcal L}
\newcommand{\cF}{\mathscr F}
\newcommand{\res}{\mathfrak R}

\begin{document}
\thispagestyle{firstpage}
\begin{flushleft}
{\small\scshape\color{muted}Probability \quad / \quad Research note}\par
\vspace{0.6em}
{\fontsize{22}{26}\selectfont\bfseries A fixed-point correction\par
for speed monotonicity}\par
\vspace{0.5em}
{\large Galton--Watson trees with at least two offspring}\par
\vspace{0.7em}
{\small\color{muted}25 September 2026 \quad\textperiodcentered\quad Working draft}
\end{flushleft}
\vspace{-.3em}
\noindent\rule{\textwidth}{0.4pt}
\vspace{.2em}

\begin{abstract}
We record an argument for strict monotonicity of the speed of a biased random
walk on a Galton--Watson tree whose offspring distribution has finite mean and
is supported on $\{2,3,\ldots\}$. The main step is an inverse-square reweighting
inequality for the escape-probability fixed point. The factor two comes from
retaining the two-child structure under reweighting, rather than from a lower
bound on the escape probability. A symmetric speed formula then suggests an
explicit correction that vanishes on every fixed point. This correction makes
the distributional variation positive and the explicit parameter variation
negative. Comparing nearby fixed points yields a strictly negative upper
right Dini derivative throughout $1<\lambda<m$, without differentiating the
infinite-depth escape probability. A known low-parameter theorem completes the
monotonicity statement. Classical differentiability is invoked separately,
under an exponential-moment assumption.
\end{abstract}

\noindent{\footnotesize\color{muted}\emph{Status.} This working note records an AI-assisted proof argument. The proof and its novelty have not undergone
independent review; the cited results are distinguished from the argument
presented here.}

\section{Statement and proof architecture}\label{sec:statement}

Let $\mathbb T$ be a Galton--Watson tree with offspring law $(p_k)_{k\ge0}$.
Throughout, assume
\begin{equation}\label{eq:assumptions}
 \PP(\nu\ge2)=1,\qquad m:=\E\nu<\infty.
\end{equation}
At a vertex with $k$ children, the $\lambda$-biased walk jumps to its parent
with probability $\lambda/(\lambda+k)$ and to each child with probability
$1/(\lambda+k)$. Thus increasing $\lambda$ increases the bias towards the root.
The choice of reflection convention at the root does not affect the speed.
For $0<\lambda<m$, write
\begin{equation}\label{eq:speed-def}
 v(\lambda)=\lim_{n\to\infty}\frac{|W_n|}{n}.
\end{equation}
The existence, determinism and positivity of this speed are part of the
standard Galton--Watson random-walk theory recalled in
\cite[Section~1]{Aidekon}. Set $v(0)=1$.

For a real-valued function $f$, use the upper right Dini derivative
\begin{equation}\label{eq:dini-def}
 \Dplus f(\lambda)
 :=\limsup_{\varepsilon\downarrow0}
       \frac{f(\lambda+\varepsilon)-f(\lambda)}{\varepsilon}.
\end{equation}

\begin{theorem}\label{thm:main}
Under \eqref{eq:assumptions}, for every $1<\lambda<m$ there exists
$\delta_\lambda>0$ such that
\begin{equation}\label{eq:main-dini}
 \Dplus v(\lambda)\le-\delta_\lambda<0.
\end{equation}
Consequently, $v$ is strictly decreasing on $[0,m)$.
\end{theorem}

\begin{corollary}\label{cor:derivative}
If, in addition, $\E e^{\theta\nu}<\infty$ for some $\theta>0$, then
$v$ is differentiable on $(0,m)$ and
\begin{equation}\label{eq:negative-derivative}
 v'(\lambda)<0,\qquad 0<\lambda<m.
\end{equation}
\end{corollary}

The additional moment assumption in \cref{cor:derivative} is used only for the
classical differentiability theorem of Bowditch--Tokushige
\cite[Theorem~1.1]{BT}. It is not used in the fixed-point inequalities or in
\eqref{eq:main-dini}.

\subsection*{The two mechanisms to keep in view}

Let $\beta_\lambda$ be the probability, on a fixed tree, of never hitting the
parent when the walk starts at the root. Fix $\lambda>1$ and put $c=\lambda-1$.
The probabilistic core of the proof is
\begin{equation}\label{eq:roadmap-tilt}
 \boxed{\quad
 \E\beta_\lambda
 <2\,
 \frac{\E[\beta_\lambda/(c+\beta_\lambda)^2]}
      {\E[1/(c+\beta_\lambda)^2]}.
 \quad}
\end{equation}
Its proof preserves the two-child structure under a change of measure; it
does not estimate $\partial_\lambda\beta_\lambda$.

The analytic core is to alter the speed-comparison functional by a quantity
that is \emph{zero at every recursive fixed point}. For an explicitly chosen
smooth function $u$, this quantity is
\begin{equation}\label{eq:roadmap-residual}
 \res_u(\alpha,\rho)
 =\E_\rho u(X)
  -\E_{\nu,\rho}u\biggl(
      \frac{\sum_{i=1}^{\nu}X_i}{\alpha+\sum_{i=1}^{\nu}X_i}
    \biggr).
\end{equation}
Subtracting $2\res_u$ leaves the value on the true fixed-point curve unchanged,
but changes both off-curve derivatives in a useful way. The choice
\begin{equation}\label{eq:roadmap-corrector}
 u'(x)=\E\frac{\beta_\lambda}
 {(\lambda-1+x+\beta_\lambda-x\beta_\lambda)^2}
\end{equation}
produces an exact cancellation. The same inequality
\eqref{eq:roadmap-tilt} then controls \emph{both} remaining signs.

The order of the proof is therefore: symmetric speed representation,
two-child reweighting, explicit correction, and finally a finite-parameter
comparison of nearby fixed points. No assertion that a root-only response is
negative is needed.

\section{Escape probabilities and a symmetric speed formula}\label{sec:reduction}

Adjoin an artificial parent to the root $e$. Let $\beta_\lambda(T)$ denote the
quenched probability of never visiting that parent, starting from $e$.
For $0<\lambda<m$, $0<\beta_\lambda<1$ almost surely. With
$\pi_\lambda=\Law(\beta_\lambda)$, the recursive distributional equation is
\begin{equation}\label{eq:rde}
 X\ \overset{\mathrm d}=\ F_\lambda(S),\qquad
 F_\lambda(s)=\frac{s}{\lambda+s},\qquad
 S=\sum_{i=1}^{\nu}X_i,
\end{equation}
where $X_i$ are independent with law $\pi_\lambda$ and independent of $\nu$.
The RDE follows by a first-step decomposition, including repeated returns
from children to the root.

For $\alpha>1$ and a probability measure $\rho$ on $[0,1]$, define
\begin{align}
 t(x,y)&=x+y-xy,
 & A_\alpha(x,y)&=\alpha-1+t(x,y),\label{eq:kernels}\\
 N_\alpha(\rho)&=\E_\rho\frac{XY}{A_\alpha(X,Y)},
 & Z_\alpha(\rho)&=\E_\rho\frac{t(X,Y)}{A_\alpha(X,Y)},\label{eq:NZ}
\end{align}
where $X,Y$ are independent with law $\rho$. If $Z_\alpha(\rho)>0$, write
$\mathcal V_\alpha(\rho)=N_\alpha(\rho)/Z_\alpha(\rho)$.

\begin{proposition}[Symmetric representation]\label{prop:symmetric}
For $1<\lambda<m$,
\begin{equation}\label{eq:symmetric-speed}
 v(\lambda)=\mathcal V_\lambda(\pi_\lambda)
 =\frac{\E[XY/A_\lambda(X,Y)]}
       {\E[t(X,Y)/A_\lambda(X,Y)]}.
\end{equation}
\end{proposition}

\begin{proof}
A\"id\'ekon's formula \cite[Theorem~1.1]{Aidekon} reads
\begin{equation}\label{eq:aidekon}
 v(\lambda)=
 \frac{\E[(\nu-\lambda)X_0/(\lambda-1+X_0+S)]}
      {\E[(\nu+\lambda)X_0/(\lambda-1+X_0+S)]},
\end{equation}
where $X_0,X_1,\ldots$ are independent copies of $\beta_\lambda$ and
$S=\sum_{i=1}^{\nu}X_i$. Conditional exchangeability gives
\begin{equation}\label{eq:exchange-speed}
 \E\frac{\nu X_0}{\lambda-1+X_0+S}
 =\E\frac{S}{\lambda-1+X_0+S}.
\end{equation}
Set $X=X_0$ and $Y=S/(\lambda+S)$. By \eqref{eq:rde}, $X,Y$ are independent
with law $\pi_\lambda$, and
\begin{equation}\label{eq:denom-change}
 \lambda-1+X+S=\frac{A_\lambda(X,Y)}{1-Y}.
\end{equation}
The numerator in \eqref{eq:aidekon} is consequently
\[
 \lambda\E\frac{Y-X+XY}{A_\lambda(X,Y)}
 =\lambda N_\lambda(\pi_\lambda),
\]
since the $Y-X$ term vanishes by symmetry. The denominator is
$\lambda Z_\lambda(\pi_\lambda)$.
\end{proof}

\begin{lemma}[Order and continuity of the fixed-point curve]\label{lem:continuity}
For $0<\alpha<\gamma<m$, escape probabilities can be coupled on the same tree
so that $\beta_\gamma\le\beta_\alpha$. Moreover,
$\lambda\mapsto\pi_\lambda$ is continuous in the $1$-Wasserstein metric.
In particular, $v$ is continuous on $(1,m)$.
\end{lemma}

\begin{proof}
For a finite-depth tree, impose escape value $1$ at the terminal generation
and iterate the map in \eqref{eq:rde} towards the root. The map is increasing
in every child value and decreasing in the parameter. Induction gives the
claimed order at every depth. The finite-depth probabilities decrease to
the probability of escaping the parent forever, giving the same order in
the limit.

The annealed escape probability $\E\beta_\lambda$ is continuous by
\cite[Proposition~2.9]{BT}. In the leafless case the proof uses
\cite[Lemma~4.7]{BT}, without an exponential-moment assumption. For ordered
random variables, the coupling just constructed gives
\begin{equation}\label{eq:W1}
 W_1(\pi_\alpha,\pi_\gamma)
 =\E\beta_\alpha-\E\beta_\gamma.
\end{equation}
Thus the fixed-point laws are $W_1$-continuous. The kernels in
\eqref{eq:NZ} are bounded and continuous locally in $\alpha>1$, and the
normalizing denominator is positive at a nonzero fixed point. This proves
the stated continuity of $v$.
\end{proof}

\section{The two-child reweighting inequality}\label{sec:tilt}

This section contains the only sign estimate that depends on the assumption
$\nu\ge2$. In particular, no maximum offspring or inverse escape-probability
moment enters the argument.

\begin{lemma}[Inverse-square reweighting]\label{lem:tilt}
Fix $\lambda>1$, put $c=\lambda-1$, and let $X\in(0,1)$ satisfy
\eqref{eq:rde} with $\nu\ge2$. Define
\begin{equation}\label{eq:mu-k}
 \mu=\E X,\qquad
 k=\frac{\E[X/(c+X)^2]}{\E[1/(c+X)^2]}.
\end{equation}
Then $\mu<2k$.
\end{lemma}

\begin{proof}
\emph{Step 1: retain two children after reweighting.}
Let $M$ be the probability measure obtained by weighting the law of $X$ with
$(c+X)^{-2}$, and write
\begin{equation}\label{eq:tilted-moments}
 k=\E_MX,\qquad j=\E_MX^2,\qquad \ell=\E_MX^3.
\end{equation}
Let $Q$ be the probability measure on $(\nu,X_1,X_2,\ldots)$ obtained by
weighting its original law with $(c+S)^{-2}$. Write
$S=X_1+H$, where $H=\sum_{i=2}^{\nu}X_i$ is independent of $X_1$.
The $X_1$-marginal under $Q$, relative to $M$, has density proportional to
\begin{equation}\label{eq:tilt-ratio}
 R(x)=(c+x)^2\E(c+x+H)^{-2}
     =\E\biggl[\biggl(\frac{c+x}{c+x+H}\biggr)^2\biggr].
\end{equation}
This is increasing in $x$. Consequently
\begin{equation}\label{eq:first-child}
 \E_QX_1=\frac{\E_M[XR(X)]}{\E_MR(X)}\ge \E_MX=k.
\end{equation}
For completeness, the covariance used here is nonnegative because, for
independent $U,V$ with law $M$,
$2\Cov_M(X,R(X))=\E[(U-V)(R(U)-R(V))]\ge0$.

Conditional exchangeability of the children now gives the decisive estimate
\begin{equation}\label{eq:two-children}
 \boxed{\quad
 \E_QS
 =\frac{\E[\nu X_1(c+S)^{-2}]}{\E[(c+S)^{-2}]}
 \ge2\E_QX_1\ge2k.
 \quad}
\end{equation}
The factor two in this inequality is the branching input of the proof.

By \eqref{eq:rde}, $S$ has the same law as $\lambda X/(1-X)$, and
$c+S$ has the same law as $(c+X)/(1-X)$. Thus
\begin{equation}\label{eq:QS}
 \E_QS=\lambda\frac{\E_M[X(1-X)]}{\E_M[(1-X)^2]}
     =\lambda\frac{k-j}{1-2k+j}.
\end{equation}
Combining \eqref{eq:two-children} and \eqref{eq:QS} yields
\begin{equation}\label{eq:moment-constraint}
 \lambda(k-j)\ge2k(1-2k+j),
 \qquad
 j\le\bar j:=\frac{k(c-1+4k)}{c+1+2k}.
\end{equation}
Since $j\ge k^2$ and
\begin{equation}\label{eq:moment-feasible}
 \bar j-k^2=\frac{k(1-k)(c+2k-1)}{c+1+2k},
\end{equation}
we necessarily have $c+2k-1\ge0$.

\emph{Step 2: undo the reweighting.}
Returning from $M$ to the original law gives
\begin{equation}\label{eq:undo-tilt}
 \mu=
 \frac{\E_M[X(c+X)^2]}{\E_M[(c+X)^2]}
 =\frac{c^2k+2cj+\ell}{c^2+2ck+j}.
\end{equation}
If $k\ge1/2$, the conclusion follows from $\mu<1$. Suppose henceforth that
$0<k<1/2$. Since $\ell\le j$,
\begin{align}
 &2k(c^2+2ck+j)-(c^2k+2cj+\ell)\notag\\
 &\hspace{1em}\ge c^2k+4ck^2+(2k-2c-1)j\notag\\
 &\hspace{1em}\ge c^2k+4ck^2+(2k-2c-1)\bar j,
 \label{eq:gap-bound}
\end{align}
where the last step uses $2k-2c-1<0$ and \eqref{eq:moment-constraint}.
Set $d=c+2k-1\ge0$. A direct rearrangement gives
\begin{align}
 &c^2k+4ck^2+(2k-2c-1)\bar j\notag\\
 &\quad=\frac{k}{c+1+2k}\bigl[
 d^3+2d^2+\{2+2k(1-2k)\}d
                  +2(1-k)(1-2k)\bigr]>0.
 \label{eq:positive-polynomial}
\end{align}
Every term in brackets is nonnegative and the last is strictly positive.
Equations \eqref{eq:undo-tilt}--\eqref{eq:positive-polynomial} imply $\mu<2k$.
\end{proof}

\begin{remark}\label{rem:why-two}
The children need not remain independent under $Q$. Independence is used
under the original measure in \eqref{eq:tilt-ratio}; conditional exchangeability
is preserved by the symmetric tilt. The lower bound
$\nu\ge2$ is needed in \eqref{eq:two-children}. Merely knowing a support bound
for $X$ does not supply either of these two pieces of information. If
single-child vertices are allowed, the same argument retains only
$\E_QS\ge \E_QX_1$ and does not prove \cref{lem:tilt}.
\end{remark}

\section{From reweighting to a positive variation kernel}\label{sec:kernel}

For the rest of the proof fix $1<\lambda<m$ and abbreviate
\begin{equation}\label{eq:basepoint}
 c=\lambda-1,\quad \pi=\pi_\lambda,\quad
 N=N_\lambda(\pi),\quad Z=Z_\lambda(\pi),\quad v=v(\lambda),
 \quad A=A_\lambda.
\end{equation}
All unqualified copies $X,Y,X_i$ have law $\pi$. Define
\begin{equation}\label{eq:J-h}
 h_v(y)=(1+v)y-v,
 \qquad
 (\cJ f)(x)=\E\frac{f(Y)}{A(x,Y)^2}.
\end{equation}

\begin{lemma}[Uniform positive margin]\label{lem:positive-kernel}
With $k$ as in \cref{lem:tilt},
\begin{equation}\label{eq:threshold-k}
 \frac{v}{1+v}<k.
\end{equation}
In particular, writing $\eta=(1+v)k-v>0$,
\begin{equation}\label{eq:J-positive}
 (\cJ h_v)(x)\ge\frac{\eta}{\lambda^2}>0,
 \qquad 0\le x\le1.
\end{equation}
\end{lemma}

\begin{proof}
Put $W=\E[X/A(X,Y)]$. Symmetry gives $N+Z=2W$, hence
$v/(1+v)=N/(2W)$. For fixed $x$, $y\mapsto A(x,y)^{-1}$ is decreasing.
The one-dimensional covariance inequality therefore gives
\[
 \E\frac{Y}{A(x,Y)}\le\mu\E\frac1{A(x,Y)},\qquad \mu=\E X.
\]
Multiplying by $x$ and averaging proves $N\le\mu W$. By \cref{lem:tilt},
\begin{equation}\label{eq:v-mu-k}
 \frac{v}{1+v}=\frac N{2W}\le\frac\mu2<k.
\end{equation}

Let $k_x=\E[Y/A(x,Y)^2]/\E[1/A(x,Y)^2]$. Relative to the measure $M$ in
\cref{lem:tilt}, this tilt has density proportional to
$\{(c+y)/A(x,y)\}^2$, an increasing function of $y$, since
\begin{equation}\label{eq:ratio-derivative}
 \frac{\partial}{\partial y}\frac{c+y}{A(x,y)}
 =\frac{\lambda x}{A(x,y)^2}\ge0.
\end{equation}
Thus $k_x\ge k$. As $A(x,y)\le\lambda$,
\[
 (\cJ h_v)(x)
 =\E[A(x,Y)^{-2}]\bigl((1+v)k_x-v\bigr)
 \ge\frac{\eta}{\lambda^2}.
\]
\end{proof}

\section{An explicit correction and two exact cancellations}\label{sec:correction}

The choice of correction is suggested by differentiating the symmetric
speed comparison. Indeed,
\begin{equation}\label{eq:variation-kernel}
 \frac{\partial}{\partial x}
       \frac{xy-vt(x,y)}{A(x,y)}
 =\frac{y^2+c h_v(y)}{A(x,y)^2}.
\end{equation}
Accordingly set
\begin{equation}\label{eq:Psi}
 \Psi=\cJ(y^2+c h_v(y)),\qquad
 b(x)=\frac{x(1-x)}{\lambda},\qquad
 \mathsf C=\E\frac{v\,t(X,Y)-XY}{A(X,Y)^2}.
\end{equation}
The quantity $\mathsf C$ is the explicit parameter derivative of
$N_\alpha(\pi)-vZ_\alpha(\pi)$ at $\alpha=\lambda$.

For a bounded function $r$, define
\begin{equation}\label{eq:L}
 (\cL r)(x)=\E\biggl[
 \frac{\nu\lambda}{(\lambda+x+S_{\nu-1})^2}
 r\biggl(\frac{x+S_{\nu-1}}{\lambda+x+S_{\nu-1}}\biggr)
 \biggr],
 \qquad S_{\nu-1}=\sum_{i=2}^{\nu}X_i.
\end{equation}
This is the operator produced by differentiating one RDE update with respect
to the input law. Its meaning will be used directly through a finite
transport formula in \cref{sec:comparison}.

Choose
\begin{equation}\label{eq:r-u}
 \boxed{\quad r(x)=\cJ\id(x)=\E\frac{Y}{A(x,Y)^2}.\quad}
\end{equation}
A convenient primitive, normalized by $u(0)=0$, is
\begin{equation}\label{eq:u-explicit}
 u(x)=\int_0^x r(s)\dd s
     =\E\frac{xY}{(c+Y)A(x,Y)}.
\end{equation}
Both $r$ and $u$ are smooth and bounded on $[0,1]$.

\begin{proposition}[Cancellation of the distributional variation]\label{prop:cancel-law}
The correction \eqref{eq:r-u} satisfies
\begin{equation}\label{eq:Lr}
 \cL r=\cJ\bigl(y(1-y)\bigr),\qquad
 (I-\cL)r=\cJ(y^2).
\end{equation}
Consequently,
\begin{equation}\label{eq:law-margin}
 \boxed{\quad
 \Psi-(I-\cL)r=c\,\cJ h_v,
 \qquad
 \inf_{x\in[0,1]}\{\Psi-(I-\cL)r\}(x)>0.
 \quad}
\end{equation}
\end{proposition}

\begin{proof}
The elementary denominator identity
\begin{equation}\label{eq:denom-F}
 A(F_\lambda(s),y)
 =\frac{\lambda(c+s+y)}{\lambda+s}
\end{equation}
cancels the derivative of $F_\lambda$ against the squared kernel. Using an
additional independent copy $Y$ in \eqref{eq:L} gives
\begin{align}
 (\cL r)(x)
 &=\frac1\lambda\E\frac{\nu Y}{(c+x+S_{\nu-1}+Y)^2}\notag\\
 &=\frac1\lambda\E\frac{S}{(c+x+S)^2}\notag\\
 &=\E\frac{Y(1-Y)}{A(x,Y)^2}.\label{eq:Lr-proof}
\end{align}
The middle equality uses exchangeability of the $\nu$ summands conditional
on $\nu$, and the last uses \eqref{eq:rde}. Equations \eqref{eq:Psi} and
\eqref{eq:J-positive} now give \eqref{eq:law-margin}.
\end{proof}

\begin{proposition}[Cancellation of the parameter variation]\label{prop:cancel-parameter}
Define
\begin{equation}\label{eq:K}
 K=2\E[b(X)r(X)]-\mathsf C.
\end{equation}
Then
\begin{align}
 \lambda K
 ={}&2\E\bigl[X(1-X)(\cJ h_v)(X)\bigr]\notag\\
 &+(1+v)\E\frac{XY(1-X)(1-Y)}{A(X,Y)^2}>0.
 \label{eq:K-positive}
\end{align}
\end{proposition}

\begin{proof}
By symmetry,
\begin{equation}\label{eq:K-expanded}
 K=\E\frac{XY\{1+(2-X-Y)/\lambda\}-v\,t(X,Y)}{A(X,Y)^2}.
\end{equation}
Since $\lambda-A(x,y)=(1-x)(1-y)$ and
$\E[(XY-vt(X,Y))/A(X,Y)]=N-vZ=0$, multiplying
\eqref{eq:K-expanded} by $\lambda$ removes the term with a single power
of $A$ in the denominator. The remaining numerator satisfies the pointwise
identity
\begin{align}
 &(1-x)(1-y)\{xy-vt(x,y)\}+xy(2-x-y)\notag\\
 &\quad=x(1-x)h_v(y)+y(1-y)h_v(x)
             +(1+v)xy(1-x)(1-y).\label{eq:pointwise-K}
\end{align}
Divide by $A(x,y)^2$ and average to obtain \eqref{eq:K-positive}.
Both terms there are positive, by \cref{lem:positive-kernel} and
$0<X,Y<1$ almost surely.
\end{proof}

\noindent\emph{The key consequence.}
The same positive kernel $\cJ h_v$ controls the two signs. There is no second
unrelated probabilistic inequality to prove. The correction cancels precisely
the $y^2$ part of \eqref{eq:variation-kernel}, and the fixed-point equality
$N-vZ=0$ turns the remaining parameter term into \eqref{eq:K-positive}.

\section{Comparing the actual fixed-point curve}\label{sec:comparison}

In this section $\lambda$, $\pi$, $v$, $r$ and $u$ are all fixed at the
base point \eqref{eq:basepoint}. In particular, the correction is \emph{not}
recomputed when $\alpha$ or $\rho$ changes.

For $\alpha>1$ and a probability measure $\rho$ on $[0,1]$, set
\begin{align}
 \res_u(\alpha,\rho)
 &=\E_\rho u(X)-\E_{\nu,\rho}u\bigl(F_\alpha(S_\rho)\bigr),
 &S_\rho&=\sum_{i=1}^{\nu}X_i,\label{eq:residual}\\
 \cF(\alpha,\rho)
 &=N_\alpha(\rho)-vZ_\alpha(\rho)-2\res_u(\alpha,\rho).
 \label{eq:corrected-functional}
\end{align}
Because the offspring law is kept fixed, every true fixed point satisfies
$\res_u(\alpha,\pi_\alpha)=0$. Hence
\begin{equation}\label{eq:on-curve}
 \boxed{\quad
 \cF(\alpha,\pi_\alpha)
 =Z_\alpha(\pi_\alpha)\{v(\alpha)-v(\lambda)\}.
 \quad}
\end{equation}
This is the reason the correction is legitimate.

\subsection{A transport formula and its local stability}

Write $\cL_{\alpha,\rho}$ for \eqref{eq:L} with parameter $\alpha$ and
child law $\rho$. Define
\begin{equation}\label{eq:G-field}
 G_{\alpha,\rho}(x)
 =\int\frac{y^2+(\alpha-1)h_v(y)}{A_\alpha(x,y)^2}\,\rho(\mathrm dy)
   -r(x)+(\cL_{\alpha,\rho}r)(x).
\end{equation}
At the base point, \cref{prop:cancel-law} states that
\begin{equation}\label{eq:G-base}
 G_{\lambda,\pi}=c\,\cJ h_v,
 \qquad \gamma:=\inf_{x\in[0,1]}G_{\lambda,\pi}(x)>0.
\end{equation}

\begin{lemma}[Transport and continuity]\label{lem:transport}
The map $(\alpha,\rho)\mapsto G_{\alpha,\rho}$ is continuous at
$(\lambda,\pi)$, with values in $C([0,1])$ equipped with the uniform norm and
with $\rho$ equipped with $W_1$.
Moreover, if $X_0\le X_1$ is a coupling, $X_s=(1-s)X_0+sX_1$, and
$\rho_s=\Law(X_s)$, then
\begin{equation}\label{eq:transport-formula}
 \frac{\mathrm d}{\mathrm ds}\cF(\alpha,\rho_s)
 =2\E\bigl[(X_1-X_0)G_{\alpha,\rho_s}(X_s)\bigr].
\end{equation}
Consequently, $\cF(\alpha,\cdot)$ is increasing along any such ordered
interpolation lying in a sufficiently small neighborhood of $\pi$, provided
$\alpha$ is sufficiently close to $\lambda$.
\end{lemma}

\begin{proof}
Fix a compact parameter interval $I\subset(1,\infty)$ containing $\lambda$,
and let $a=\inf I>1$. The double-integral kernel in \eqref{eq:G-field} is
bounded and jointly continuous on $I\times[0,1]^2$.
For the recursive term, conditional on $\nu=n$ its integrand is bounded by
\begin{equation}\label{eq:nu-bound}
 \frac{n}{a}\|r\|_\infty.
\end{equation}
For each fixed $n$ it is jointly continuous in the parameter, the external
coordinate $x$ and the $n-1$ child coordinates. Integrating against
$\rho^{\otimes(n-1)}$ therefore gives uniform convergence in $x$ when the
parameter and the law converge. One way to see uniformity is to use uniform
continuity of the integrand and then a finite net in $[0,1]$.
Truncating the sum at $n=M$ leaves a remainder bounded, uniformly in
$\alpha\in I$, $x$ and $\rho$, by
\begin{equation}\label{eq:nu-tail}
 \frac{\|r\|_\infty}{a}\E[\nu\mathbf1_{\{\nu>M\}}]\longrightarrow0.
\end{equation}
This proves the asserted continuity. Notice that only $\E\nu<\infty$ is
used here.

For the transport identity, take independent copies of the coupled pair
$(X_0,X_1)$ for all child coordinates and for the two variables in
\eqref{eq:NZ}. Differentiation of $N_\alpha-vZ_\alpha$ gives twice the first
term of \eqref{eq:G-field}. Differentiation of $\E_{\rho_s}u(X)$ gives
$\E[(X_1-X_0)r(X_s)]$. Finally, the chain rule for the RDE update produces
$\E[(X_1-X_0)(\cL_{\alpha,\rho_s}r)(X_s)]$ after conditional
exchangeability. Interchanging differentiation and expectation is justified
by \eqref{eq:nu-bound}, $0\le X_1-X_0\le1$ and $\E\nu<\infty$.
This proves \eqref{eq:transport-formula}.

By continuity and \eqref{eq:G-base}, $G_{\alpha,\rho}\ge\gamma/2$ uniformly
in $x$ in a sufficiently small neighborhood. Formula
\eqref{eq:transport-formula} is then nonnegative on every ordered
interpolation contained in that neighborhood.
\end{proof}

\subsection{The parameter derivative and the final comparison}

Freezing the distribution $\pi$ while varying the parameter gives
\begin{equation}\label{eq:parameter-residual}
 \left.\partial_\alpha\res_u(\alpha,\pi)\right|_{\alpha=\lambda}
 =\E\bigl[b(X)r(X)\bigr].
\end{equation}
Indeed, $\partial_\alpha F_\alpha(s)=-s/(\alpha+s)^2$, which at the fixed
point equals $-b(F_\lambda(s))$. The derivative is dominated since
$s/(\alpha+s)^2\le1/(4\alpha)$. From \cref{prop:cancel-parameter},
\begin{equation}\label{eq:parameter-negative}
 \left.\partial_\alpha\cF(\alpha,\pi)\right|_{\alpha=\lambda}
 =\mathsf C-2\E[br]=-K<0.
\end{equation}

\begin{proof}[Proof of \cref{thm:main} on $(1,m)$]
Let $\varepsilon>0$ be small enough that $\lambda+\varepsilon<m$.
By \cref{lem:continuity}, on the same tree
\[
 U_\varepsilon:=\beta_{\lambda+\varepsilon}\le
 V:=\beta_\lambda,
 \qquad \E(V-U_\varepsilon)\longrightarrow0.
\]
The interpolation $X_s=(1-s)U_\varepsilon+sV$ satisfies
\begin{equation}\label{eq:interpolation-ball}
 W_1(\Law(X_s),\pi)
 \le\E|X_s-V|
 =(1-s)\E(V-U_\varepsilon)
 \le\E(V-U_\varepsilon).
\end{equation}
Thus, for sufficiently small $\varepsilon$, the entire interpolation is
inside the neighborhood in \cref{lem:transport}. With
$\alpha=\lambda+\varepsilon$, that lemma gives
\begin{equation}\label{eq:law-comparison}
 \cF(\lambda+\varepsilon,\pi_{\lambda+\varepsilon})
 \le\cF(\lambda+\varepsilon,\pi).
\end{equation}
Since $\cF(\lambda,\pi)=0$, \eqref{eq:parameter-negative} yields
\[
 \cF(\lambda+\varepsilon,\pi)=-K\varepsilon+o(\varepsilon).
\]
Using \eqref{eq:on-curve}, we conclude that
\begin{equation}\label{eq:finite-comparison}
 Z_{\lambda+\varepsilon}(\pi_{\lambda+\varepsilon})
 \{v(\lambda+\varepsilon)-v(\lambda)\}
 \le-K\varepsilon+o(\varepsilon).
\end{equation}
The normalizing factor converges to $Z>0$. Therefore
\begin{equation}\label{eq:explicit-dini}
 \boxed{\quad \Dplus v(\lambda)\le-\frac KZ<0.\quad}
\end{equation}
This proves \eqref{eq:main-dini}, with $\delta_\lambda=K/Z$.

A continuous function with strictly negative upper right Dini derivative at
every point of an interval is strictly decreasing there. Indeed, if
$f(b)\ge f(a)$ for $a<b$, a minimum of $f$ on $[a,b]$ is attained at some
$x<b$: if $b$ is a minimizer, then $a$ is also one. At such an $x$ all
sufficiently small right difference quotients are nonnegative, a
contradiction. By \cref{lem:continuity}, this argument applies to $v$ on
$(1,m)$.
\end{proof}

\subsection{The low-parameter interval and classical derivatives}

Let $m_1=\min\{k:p_k>0\}\ge2$. Song--Wang--Xiang
\cite[Theorem~1.2 and Lemma~2.5]{SWX} establish $v'(\lambda)<0$ for
\begin{equation}\label{eq:low-range}
 0<\lambda\le a_{m_1},\qquad
 a_{m_1}:=\frac{m_1}{1+\sqrt{1-1/m_1}}>1,
\end{equation}
and strict decrease on $[0,a_{m_1}]$. Their derivative statement and its
justification are given immediately before, and in, Lemma~2.5 of the cited
preprint. This interval overlaps $(1,m)$, so its strict monotonicity and
\eqref{eq:explicit-dini} complete \cref{thm:main} on $[0,m)$.

Under the exponential-moment assumption of \cref{cor:derivative},
\cite[Theorem~1.1]{BT} gives differentiability on $(0,m)$: its lower endpoint
is zero in the leafless case. At every $\lambda>1$ the derivative equals the
Dini derivative and is negative by \eqref{eq:explicit-dini}. The remaining
parameters are covered by \eqref{eq:low-range}. This proves
\cref{cor:derivative}.

\section{What the correction accomplishes}\label{sec:meaning}

\paragraph{It retains the full recursive response.}
At a purely formal level, writing $D=-\partial_\lambda\beta_\lambda$ and
differentiating \eqref{eq:rde} gives
\begin{equation}\label{eq:formal-response}
 D=\frac{\beta(1-\beta)}{\lambda}
    +\frac{(1-\beta)^2}{\lambda}\sum_{i=1}^{\nu}D_i.
\end{equation}
Keeping only the first term discards the change in every descendant
conductance. The correction does not discard that change. Instead,
$\res_u(\alpha,\pi_\alpha)=0$ supplies an exact identity along the full
fixed-point curve. The proof above uses finite transport and parameter
variations of bounded kernels, so \eqref{eq:formal-response} is explanatory
only, not an assumption or an analytic step in the proof.

\paragraph{The offspring assumptions have separate roles.}
The inequality $\nu\ge2$ enters precisely in \eqref{eq:two-children}.
The finite-mean assumption permits \eqref{eq:nu-tail}, the transport
interchange, and the application of the speed formula. An exponential
moment is used only to obtain a classical derivative everywhere via the
external differentiability theorem. No bounded-offspring assumption is
made, and the argument does not address laws with $p_1>0$.

\paragraph{The proof depends on one probabilistic inequality, not two.}
The real sign input is \cref{lem:tilt}. It gives $\cJ h_v>0$; the first
cancellation \eqref{eq:law-margin} and the second cancellation
\eqref{eq:K-positive} then give the two off-curve signs from that same input.
The entire mechanism can be read as
\[
 \begin{gathered}
 \text{two children under a tilt}
 \ \Longrightarrow\ \mu<2k
 \ \Longrightarrow\ \cJ h_v>0,\\[0.4em]
 \cJ h_v>0\ \Longrightarrow\
 \begin{cases}
   \text{positive distributional variation},\\
   \text{negative explicit parameter variation}.
 \end{cases}
 \end{gathered}
\]
The vanishing of the residual on fixed points transfers these off-curve
signs back to the actual speed.

\begin{remark}[Regular-tree check]\label{rem:regular}
For the deterministic $d$-ary tree, $\beta_\lambda=1-\lambda/d$ and
\eqref{eq:symmetric-speed} gives
\[
 v(\lambda)=\frac{d-\lambda}{d+\lambda},\qquad
 v'(\lambda)=-\frac{2d}{(d+\lambda)^2}.
\]
The reweighted mean $k$ equals $\E\beta_\lambda$ in this case, so the strict
inequality of \cref{lem:tilt} is consistent even at a point-mass fixed point.
\end{remark}

\appendix
\section{Algebraic verification and reference conventions}\label{sec:appendix}

The accompanying file \texttt{check\_algebra.py} verifies the rational and
polynomial identities used in the argument by exact symbolic arithmetic.
In particular, it checks the transport derivative
\eqref{eq:variation-kernel}, the feasible-moment identity
\eqref{eq:moment-feasible}, the positive expansion
\eqref{eq:positive-polynomial}, the denominator cancellation
\eqref{eq:denom-F}, and the numerator identity \eqref{eq:pointwise-K}.
It is a check on algebra only: the independence, reweighting, continuity and
finite-parameter comparison are proved in the text, not certified by that
script.

All numbered references to external theorems and lemmas refer to the exact
arXiv versions specified below. This avoids ambiguity between the numbering
or titles of a preprint and its published version. In particular, the
escape-probability continuity input and the differentiability input from
\cite{BT} have different moment requirements and are used separately.

\begingroup
\small
\begin{thebibliography}{BT20}
\setlength{\itemsep}{0.5em}

\bibitem[A\"id14]{Aidekon}
\'{E}.~A\"id\'ekon.
\emph{Speed of the biased random walk on a Galton--Watson tree}.
Probability Theory and Related Fields \textbf{159} (2014), 597--617.
\href{https://doi.org/10.1007/s00440-013-0515-y}{doi:10.1007/s00440-013-0515-y}.
Theorem numbering used here:
\href{https://arxiv.org/abs/1111.4313v4}{arXiv:1111.4313v4} (18 January 2013).

\bibitem[BT20]{BT}
A.~Bowditch and Y.~Tokushige.
\emph{Differentiability of the speed of biased random walks on
Galton--Watson trees}.
\href{https://arxiv.org/abs/1906.07913v2}{arXiv:1906.07913v2}
(1 May 2020).
Theorem~1.1, Proposition~2.9 and Lemma~4.7.

\bibitem[SWX16]{SWX}
H.~Song, L.~Wang and K.~Xiang.
\emph{Monotonicity of the speed for biased random walk on
Galton--Watson tree}.
\href{https://arxiv.org/abs/1610.08151v1}{arXiv:1610.08151v1}
(26 October 2016).
Theorem~1.2 and Lemma~2.5.
Published under the title \emph{The speed of a biased walk on a
Galton--Watson tree without leaves is monotonic for low values of bias},
Journal of Applied Probability \textbf{62} (2025), no.~3;
\href{https://doi.org/10.1017/jpr.2024.113}{doi:10.1017/jpr.2024.113}.

\end{thebibliography}
\endgroup

\end{document}
